% Computer Science A --- Git version control basics (\input from \texttt{main.tex}).

\runningheader{Computer Science A}{Git version control}{Page \thepage\ of \numpages}

\uplevel{\par\textbf{Source control and Git (free response)}\par}

\question[4]
What is source control management? Why is it useful?
\makeemptybox{1.25in}
\begin{solution}
  Systems (e.g.\ Git) that record history of changes, enable collaboration, branches, review, rollback, and traceability---reducing loss of work and coordinating teams.
\end{solution}

\question[4]
What are the main sections of a Git project?
\makeemptybox{1.25in}
\begin{solution}
  Working tree (edited files), staging area (index), local repository (\texttt{.git} commits), and remotes (shared repos such as origin); data moves via add/commit/push/pull.
\end{solution}

\question[4]
What are some common operations you would be performing when using Git?
\makeemptybox{1.45in}
\begin{solution}
  Clone, pull/fetch, branch/checkout, status, diff, add, commit, merge or rebase, push, tag, and resolve conflicts---daily workflow for integrating changes safely.
\end{solution}

\newpage

\question[4]
What are some common Linux commands? Why are they useful?
\makeemptybox{1.45in}
\begin{solution}
  Examples: \texttt{ls}, \texttt{cd}, \texttt{pwd}, \texttt{cp}, \texttt{mv}, \texttt{rm}, \texttt{mkdir}, \texttt{grep}, \texttt{chmod}, \texttt{ssh}. They navigate filesystems, search logs, set permissions, and automate server/CI tasks.
\end{solution}

\question[4]
The command to create a new directory in Linux is:
\begin{choices}
  \CorrectChoice \texttt{mkdir}
  \choice \texttt{touch}
  \choice \texttt{nano}
  \choice \texttt{vi}
  \choice \texttt{rm}
\end{choices}
\begin{solution}
  \texttt{mkdir} creates directories; \texttt{touch} updates/creates empty files; \texttt{nano}/\texttt{vi} are editors; \texttt{rm} removes files.
\end{solution}

\question[4]
To copy a file in Linux, you use the command:
\begin{choices}
  \CorrectChoice \texttt{cp}
  \choice \texttt{mv}
  \choice \texttt{rm}
  \choice \texttt{cat}
  \choice \texttt{mkdir}
\end{choices}
\begin{solution}
  \texttt{cp} copies; \texttt{mv} moves/renames; \texttt{rm} deletes; \texttt{cat} prints file contents; \texttt{mkdir} makes directories.
\end{solution}

\newpage

\begingradingrange{csagit}

\uplevel{\par\textbf{Git}\par}

\question[4]
What is the \texttt{git push} command used for?
\begin{choices}
  \choice To add all the files to the given folder
  \choice It is used to create a new empty repository or directory consisting of files with the hidden directory
  \CorrectChoice It is used to upload local repository content to a remote repository
  \choice It places all of the revision information in one place
\end{choices}
\begin{solution}
  \texttt{git push} sends commits from your local branch to the configured \textbf{remote} (e.g.\ \texttt{origin}).
\end{solution}

\question[4]
A developer working on the same project as you just made an official commit to the remote repository. What command gives you the updated changes?
\begin{choices}
  \CorrectChoice \texttt{git pull}
  \choice \texttt{git log}
  \choice \texttt{git push}
  \choice \texttt{git sync}
\end{choices}
\begin{solution}
  \texttt{git pull} typically fetches remote updates and integrates them into your current branch (\texttt{fetch} + \texttt{merge} or \texttt{rebase}, depending on configuration). \texttt{git sync} is not a standard Git command.
\end{solution}

\question[4]
\blank[7em] is the practice of tracking and managing changes to software code. Git is one of the most popular tools used to achieve this.
\begin{choices}
  \choice Revision tracking
  \choice Code review
  \choice Project control
  \CorrectChoice Version control
\end{choices}
\begin{solution}
  \textbf{Version control} records changes over time so teams can collaborate, review history, and revert when needed.
\end{solution}

\question[4]
What is a Git repository?
\begin{choices}
  \choice A project pipeline that helps automate the deployment process of applications
  \CorrectChoice A directory that contains a \texttt{.git} folder which tracks changes made to files within that directory
  \choice A place to plan out the requirements of specific projects between different teams working on a client project
  \choice A communication platform for developers to talk to each other about code and build projects that allow the client to see the code
\end{choices}
\begin{solution}
  The \texttt{.git} directory stores objects, refs, and config for that project; the working tree holds your files.
\end{solution}

\question[4]
What is Git?
\begin{choices}
  \choice A website
  \CorrectChoice A version control system
  \choice An operating system
  \choice Programming language
\end{choices}
\begin{solution}
  Git is a \textbf{distributed version control system} for tracking changes to files (especially source code). GitHub is a hosting service that uses Git, not Git itself.
\end{solution}

\question[4]
Which of the following statements best describes Git?
\begin{choices}
  \CorrectChoice Git is version control software that allows developers to manage changes to source code over time
  \choice Git is a tool used by development teams to write code anywhere in the world
  \choice Git is another name for GitHub, an online repository tool
  \choice Git is a database used to store data for web applications
\end{choices}
\begin{solution}
  Git is a \textbf{distributed} version control system. GitHub hosts Git repositories but is not the same as Git.
\end{solution}

\question[4]
Which of the following does \emph{not} describe git?
\begin{choices}
  \choice It is a version control software for source code
  \choice Unlike Subversion or Mercurial, git is designed to be distributed instead of centralized
  \CorrectChoice Another name for git is GitHub
  \choice The units of work that are tracked in git are called commits
\end{choices}
\begin{solution}
  Git is the tool; \textbf{GitHub} (GitLab, Bitbucket, etc.) are services built around Git repositories.
\end{solution}

\question[4]
Which of the following best describes the command \texttt{git commit}?
\begin{choices}
  \choice It pulls the changes from the remote repository to the local repository
  \choice It stages changes made in the working directory for the next commit
  \choice It creates a new branch in your repository
  \CorrectChoice It saves the staged changes to the local repository
\end{choices}
\begin{solution}
  \texttt{git commit} records a snapshot of the \textbf{index} (staged changes) in the local object database; use \texttt{git add} first to stage.
\end{solution}


\newpage


\question[4]
Which of the following best describes a merge conflict?
\begin{choices}
  \CorrectChoice It is a situation where two branches have made changes to the same part of a file and Git does not know which changes to keep
  \choice It is when Git fails to merge two branches because one is older than the other
  \choice It is the conflict that occurs when two users try to push to the same remote repository at the same time
  \choice It is a state in Git where a file is marked for deletion but is still present in the staging area
\end{choices}
\begin{solution}
  Conflicts appear when overlapping edits must be reconciled; you resolve markers in the working tree, then stage and commit.
\end{solution}

\question[4]
What are the appropriate steps to update unstaged changes to a local commit and send those changes to a remote repository?
\begin{choices}
  \choice \texttt{git add .} then \texttt{git commit -m "new commit"}
  \CorrectChoice \texttt{git add .}, then \texttt{git commit -m "new commit"}, then \texttt{git push}
  \choice \texttt{git commit -m "new commit"} then \texttt{git push}
  \choice \texttt{git commit -m "new commit"} then \texttt{git add .} then \texttt{git push}
\end{choices}
\begin{solution}
  Stage (\texttt{add}), commit locally, then \texttt{push} to publish. Committing with nothing staged (or wrong order) does not match this workflow.
\end{solution}

\question[4]
Which of the following best describes the Git workflow?
\begin{choices}
  \choice \texttt{git clone}, make changes to files, \texttt{git commit}, \texttt{git add}, \texttt{git push}
  \choice \texttt{git clone}, \texttt{git add}, make changes to files, \texttt{git push}, \texttt{git commit}
  \CorrectChoice \texttt{git clone}, make changes to files, \texttt{git add}, \texttt{git commit}, \texttt{git push}
  \choice \texttt{git add}, \texttt{git commit}, make changes to files, \texttt{git push}, \texttt{git clone}
\end{choices}
\begin{solution}
  Clone a repo, edit files, \textbf{stage} changes, \textbf{commit} locally, then \textbf{push} to share.
\end{solution}

\question[4]
What does the command \texttt{git add} do?
\begin{choices}
  \choice Creates a new repository in the current directory
  \choice Deletes the changes in your working directory
  \CorrectChoice It stages changes made in the working directory for the next commit
  \choice It connects your local repository to a remote repository
\end{choices}
\begin{solution}
  \texttt{git add} updates the \textbf{index} (staging area) so the next \texttt{git commit} includes those paths.
\end{solution}

\newpage


\question[4]
Git Bash is an application for Microsoft Windows environments which provides an emulation of a bash shell and has the Git CLI preinstalled.
\begin{choices}
  \CorrectChoice TRUE
  \choice FALSE
\end{choices}
\begin{solution}
  Git for Windows bundles \textbf{Git Bash}, a shell environment for running \texttt{git} and Unix-style commands on Windows.
\end{solution}

\question[4]
Git is a centralized version control system.
\begin{choices}
  \choice TRUE
  \CorrectChoice FALSE
\end{choices}
\begin{solution}
  Git is a \textbf{distributed} version control system (DVCS): every clone carries a full copy of history; remotes add collaboration, but there is no single required central server for the tool itself.
\end{solution}

\newpage

\uplevel{\par\textbf{Branches and remotes (practice)}\par}

\question[4]
Identify the correct series of commands to:
\begin{enumerate}
  \item switch to the \texttt{main} branch;
  \item inspect the working tree;
  \item stage and commit a change; and
  \item update the remote repository.
\end{enumerate}
\begin{choices}
  \choice \texttt{git pull}; \texttt{git add}; \texttt{git commit}; \texttt{git push}
  \choice \texttt{git edit}; \texttt{git add}; \texttt{git checkout main}; \texttt{git commit}
  \CorrectChoice \texttt{git checkout main}; \texttt{git status}; \texttt{git add}; \texttt{git commit}; \texttt{git push}
  \choice \texttt{git commit main}; \texttt{git status}; \texttt{git add}; \texttt{git push}
\end{choices}
\begin{solution}
  \texttt{checkout} (or modern \texttt{git switch}) selects the branch; \texttt{status} shows changes; \texttt{add} stages; \texttt{git commit} records a snapshot; \texttt{git push} publishes commits to the remote. (Teams often \texttt{pull} first, but the stem asks for this sequence.)
\end{solution}

\endgradingrange{csagit}
